\chapter*{Abstract}
\addcontentsline{toc}{chapter}{Abstract}

\begin{foldable}
  Knowledge distillation (KD) has emerged as a powerful technique for compressing large neural networks into compact, deployable models by training a student to mimic a teacher's outputs.
  Standard KD methods, however, assume that the practitioner has access to the full training dataset and the internal parameters of the teacher network.
  Yet, this assumption routinely fail in real-world deployments where data is proprietary, privacy-sensitive, or simply unavailable, and where teachers are exposed only as black-box APIs.
\end{foldable}

\begin{foldable}
  This thesis addresses KD under three increasingly restrictive practical constraints: (i) few-shot access to real data with a black-box teacher, (ii) complete absence of real data with a black-box teacher, and (iii) dataset distillation for text, where the goal is to synthesise a tiny surrogate dataset that preserves model training utility.
  For the first setting, we propose \emph{DivBFKD}, a framework that improves the diversity of the limited proxy data through explicit diversity regularization, enabling more effective black-box \ac{KD} from only a handful of samples per class.
  For the second setting, we propose \emph{DIPKD}, a data-free black-box \ac{KD} method that synthesises diverse training images by incorporating structured image priors, overcoming the mode collapse and low fidelity that plague existing generator-based approaches.
  For the third and final setting, we propose \emph{TAKE}, a trajectory-aware knowledge estimation method for text dataset distillation that leverages training dynamics to better preserve the knowledge encoded in large text corpora.
  Across all three settings, the proposed methods achieve state-of-the-art performance on standard benchmarks, demonstrating that effective KD is achievable even under severe practical constraints.
\end{foldable}
